%%%%%%%%%%%%%%%%%%%%%%%%%%%%%%%%%%%%%%%%%%%%%%%%%%%%%%%%%%%%%%%%%%%%%%
%  月度总结 · 阶段二（研究推进）  ——  陈思文 · 微通道液冷设计
%
%  编译：xelatex 文件名.tex   （跑两遍）
%     Overleaf：Menu -> Compiler -> XeLaTeX
%%%%%%%%%%%%%%%%%%%%%%%%%%%%%%%%%%%%%%%%%%%%%%%%%%%%%%%%%%%%%%%%%%%%%%
\documentclass[stage2]{ydzjreport}

% ===================== 只改这里：你的信息 =====================
\stagetwosetup
  {阶段二 · 研究推进}%          1 报告阶段
  {陈思文}%                   2 姓名
  {微通道液冷设计}%             3 方向（2-6 个汉字）
  {2026-08-01 至 2026-09-30}%   4 报告周期
  {热流密度分区；翼型肋片；流热耦合；网格无关性}%   5 关键词
  {完成不同热流密度分区的变截面翼型流热耦合仿真，并配合团队完成了网格无关性检验。}%     6 一句话摘要
% ============================================================

\begin{document}
\makestagetwohead

% ==================================================================
\section{本月研究任务}

\subsection{任务目标}

% 月初定的目标，2--4 条。
\begin{enumerate}
  \item 针对三种不同热流密度区域（低、高、浮动），确立变截面翼型肋片的选取思路。
  \item 仿真不同翼型构型在低流速与高流速工况下的散热差异，评估各构型的适用性。
  \item 对六种不同构型开展网格无关性检验，保证构型间对比数据的可靠性。
\end{enumerate}

\subsection{完成情况}

% 状态列只写三种：已完成 / 部分完成 / 未开始。
\begin{table}[htbp]
\centering
\caption{本月任务完成情况}
\begin{tabularx}{\textwidth}{@{}Y c Y@{}}
\toprule
\rowcolor{ydzjlight} 任务 & 状态 & 说明 \\
\midrule
三种热流密度区域的翼型肋片选取思路 & 已完成 & 形成低、高、浮动热流密度分区的差异化肋片布置方案 \\
不同翼型低/高流速散热仿真 & 已完成 & 完成多构型流热耦合仿真并输出对比数据 \\
六种构型网格无关性检验 & 已完成 & 无关性检验结果总结由队友完成，本人负责部分构型的仿真与数据积累 \\
\bottomrule
\end{tabularx}
\end{table}


% ==================================================================
\section{进阶阅读与理论支撑}

\subsection{本阶段文献清单}

% 与阶段一不同：这里每篇文献必须对应「解决了我哪个问题」。
\begin{table}[htbp]
\centering
\caption{本阶段文献}
\begin{tabularx}{\textwidth}{@{}c Y Y@{}}
\toprule
\rowcolor{ydzjlight} 序号 & 文献 & 解决了我哪个问题 \\
\midrule
1 & 《微通道液冷场景下的翼型参数》（聚焦 NACA 翼型数据） & 明确了相对厚度、相对弯度、攻角等翼型基础参数对液冷性能的影响机理，为扰流单元选型提供了参数依据 \\
\bottomrule
\end{tabularx}
\end{table}

\subsection{对本方向的启示}

通过研读 NACA 翼型数据在液冷场景下的适用性，明确了相对厚度、相对弯度与攻角是决定扰流柱换热与流阻特性的三个基础参数。据此，本月将二维翼型剖面在 SolidWorks 中拓展为直棱柱形、锥形、凹坑形三类三维扰流结构：三类结构均以周期性破坏热边界层为目标，但扰动强度与流阻代价各不相同，正好对应低、高、浮动三种热流密度区域的差异化需求。该思路使结构选型第一次有了明确的参数化抓手，也为后续的单参数扫掠研究奠定了基础。


% ==================================================================
\section{研究进展}

\subsection{已完成的工作}

% 具体做了什么，写到别人可以照着做的程度。
\begin{enumerate}
  \item 在芯片底部中央的浮动热流密度区布置变截面翼型肋片，完成分区差异化流道结构的建模与网格划分。
  \item 在低流速与高流速两种工况下，完成了相同热流密度底板下多种翼型构型的流热耦合仿真，输出温度场、压力场及各项评价指标。
  \item 输出多组不同网格数量下的测试数据，供团队开展六种构型的网格无关性检验使用。
\end{enumerate}

\subsection{关键结果}

在基准工况（热流密度 $4.0\times10^{6}$ W/m$^2$，质量流率 0.004 kg/s，入口 $x$ 向速度 1.34 m/s）下，各翼型构型（5 mm 弦长、5 个 $\times$ 6 通道布置）的仿真关键数据汇总如下表。

\begin{table}[htbp]
\centering
\caption{不同翼型构型流热耦合仿真关键结果（基准工况）}
\begin{tabularx}{\textwidth}{@{}l c c c c c@{}}
\toprule
\rowcolor{ydzjlight} 构型 & 最高温度 (K) & 底板平均温度 (K) & $h$ (W/(m$^2\cdot$K)) & $Nu$ & $PEC$ \\
\midrule
NACA0015 & 357.9 & 356.1 & 756.6 & 1196.6 & 1.00 \\
NACA0021 & 355.7 & 354.0 & 765.3 & 1211.5 & 1.10 \\
NACA0024 & 354.9 & 353.2 & 798.5 & 1264.8 & 1.19 \\
NACA4424 & 354.4 & 352.7 & 789.8 & 1251.3 & 1.21 \\
NACA6424 & 358.1 & 356.3 & 750.5 & 1186.7 & 1.00 \\
NACA6624 & 357.9 & 356.1 & 750.2 & 1186.4 & 1.01 \\
\bottomrule
\end{tabularx}
\end{table}

表中 $h$ 为对流换热系数，$Nu$ 为努塞尔数，$PEC$ 为热性能综合评价指标（以 NACA0015 为基准归一化）。此外，在高质量流率工况（1.2 倍基准质量流率）下对 NACA6624 与 NACA4424 的测试显示，换热能力虽进一步提升，但进出口压差与泵耗急剧增大，逼近压降设计红线。

\subsection{与预期的差异}

\begin{enumerate}
  \item 厚度增益符合预期：增大相对厚度（NACA0015 $\rightarrow$ NACA0024）后，最高温度由 357.9 K 降至 354.9 K，$Nu$ 提高约 5.7\%，$PEC$ 达 1.19，说明厚翼型的换热强化收益大于其摩擦损失代价。
  \item 弯度效应与预期存在偏差：预想弯度越大换热越好，但 NACA64 系列（NACA6424、NACA6624）的 $PEC$ 反而回落至 1.00 附近，最高温度回升至约 358 K。结合流场分析，高弯度翼型尾缘附近形成流动死区，分离使有效换热面积减小，抵消了弯度带来的增益。
  \item 高流速下压降激增：在 1.2 倍基准质量流率工况下，进出口压差由低流速工况的数百 Pa 量级跃升至 $10^{7}$ Pa 量级，泵耗大幅上升，说明单纯提高流速换换热的路线已逼近压降红线，能效比明显下降。
  \item 网格无关性：多组网格数量下的测试数据对比表明，网格加密到一定数量后，温度峰值与压降结果的变化幅度已很小，说明当前网格基准下解已收敛，网格划分差异不再构成结果干扰。
\end{enumerate}


% ==================================================================
\section{问题与对策}

\subsection{遇到的主要问题}

% 至少 2 条，写清卡在哪。
\begin{enumerate}
  \item 整体微通道流道内扰流柱数量较多，各翼型参数（厚度、弯度、攻角）与流动干扰相互耦合，无法从整体仿真中精准剥离单一变量对换热效率的独立影响。
  \item 构型间对比结果受网格划分差异与流动死区等因素干扰，直接下结论存在可信度风险。
\end{enumerate}

\subsection{已尝试的方案}

% 试过什么、为什么不行。与上面的问题一一对应。
\begin{enumerate}
  \item 针对问题2，配合团队完成六种构型的网格无关性检验，通过多组网格数量下的数据对比排除了网格划分对结果的干扰，确认当前数据具备构型间横向对比资格。
  \item 针对问题1，尝试在整体流道层面对比不同构型，但因多柱干扰无法解耦，宏观对比难以定位单参数主因，该路线宣告受限。
\end{enumerate}

\subsection{下一步对策}

% 打算怎么办。与上面的问题一一对应。
\begin{enumerate}
  \item 缩小研究范围，聚焦到“单个扰流柱”，建立单柱局部流热耦合模型，逐一控制变量考察厚度、弯度、攻角的独立影响，从根本上解决参数解耦问题。
  \item 在单柱规律明确后，以 $PEC$ 与压降红线为双重约束，反向完成整体通道结构的定型。
\end{enumerate}


% ==================================================================
\section{阶段性结论}

\subsection{可以确定的}

% 有证据支撑的结论，至少 2 条。
\begin{enumerate}
  \item 对热流密度进行微分并分区采用差异化翼型肋片，有效改善了局部过热与无热源区过度冷却并存的问题，底板温度分布明显优于均匀布置方案。
  \item 经网格无关性检验，确立了当前仿真边界条件下的无关性网格划分基准，后续仿真可沿用该基准保证数据可比性。
\end{enumerate}

\subsection{尚不能确定的}

% ★ 本节不可为空，至少 1 条。写清「缺什么证据才能确定」。
\begin{enumerate}
  \item 单个翼型扰流柱的温降规律尚未摸清，需要建立单柱基准模型并取得温降基准数据后方可确定。
  \item 各翼型参数（相对厚度、相对弯度）对散热效率的独立影响权重尚未定量，目前缺乏单参数扫掠的仿真证据。
\end{enumerate}


% ==================================================================
\section{下月计划}

\begin{table}[htbp]
\centering
\caption{下月工作计划}
\begin{tabularx}{\textwidth}{@{}Y Y c@{}}
\toprule
\rowcolor{ydzjlight} 事项 & 预期产出 & 时间 \\
\midrule
研究单个扰流柱的温降基准模型 & 单柱温降基准数据 & 第 1 周 \\
确定翼型各项参数（厚度、弯度、攻角）对独立散热效率的影响 & 参数影响特性曲线与数据表 & 第 2--3 周 \\
基于单柱规律完成最终通道结构的定型 & 整体通道结构设计定稿 & 第 4 周 \\
\bottomrule
\end{tabularx}
\end{table}


% ==================================================================
\begin{thebibliography}{9}
\bibitem{ref1} 微通道液冷场景下的翼型参数[R]. 2026.
\end{thebibliography}

\end{document}
